% Budget: 4--5 pp.  Function: pose the question, tell the arc end-to-end,
% state the contributions.  This is the ONLY place the chronological arc
% is told; later chapters stay deliverable-separated (D1--D4).
\chapter{Introduction}
\label{ch:introduction}

\section{Motivation}
\label{sec:motivation}
\todo{Pretrained generative models as world-model priors. src: 10_now/positioning}

\section{Problem statement}
\label{sec:problem}
\todo{Frame as WHEN AND WHY adaptation works, not CAN IT. See rubric/01-originality.}

\section{The arc}
\label{sec:arc}
\todo{AVID works -> planning -> too slow -> shortcut -> curvature -> flow/Wan
-> the pathway result -> the economics bound. src: writing/thesis-storyline}

\section{Contributions}
\label{sec:contributions}

Adapting a frozen video diffusion model into an action-conditioned world
model is not a gap in the literature: it is the contribution of
AVID~\citep{\nv{rigter2024avid}}, which trains an action-conditioned adapter
on a frozen base using only that base's predictions. We take that result as
our starting point. What is not established is \emph{when} the approach
works, \emph{why} it stops working, and whether the same composition can
carry a second conditioning variable (the sampling step size) that
would make the resulting world model fast enough to plan through. Those are
the questions this thesis answers.

\paragraph{A design principle for conditioning a frozen generative prior.}
Our central finding is architectural, and we state it as a design rule
rather than as a property of one system:

\begin{quote}
\emph{Conditioning must enter the network scale-free relative to the
residual stream.} Supplied as a per-frame modulation of \emph{normalised}
activations, an action signal survives and is used. Supplied as
cross-attention tokens, which add into an unnormalised residual stream,
the same signal survives the attention layers and is then attenuated to
irrelevance at the residual addition, at every scale and temporal alignment
we tested.
\end{quote}

We establish this causally rather than by correlation: a single-variable A/B
in which the two arms are matched on adapter contribution and on the
composition mask, so the injection pathway is the only difference
(\cref{sec:res-pathway}).

\emph{Scope of validity.} The principle is demonstrated on two backbone
families (a U-Net diffusion model and a rectified-flow DiT), at two
injection sites, for one conditioning modality (robot actions), on video
prediction. We do not test it on other modalities, on architectures without
a normalisation-modulation site, or at larger scales, and we do not claim it
transfers there without evidence.

\paragraph{A bound that architecture cannot remove.} The failures above are
mis-calibration and are fixable. Underneath them sits one that is not: under
a teacher-forced denoising objective the actions account for a very small
fraction of the loss on this data, so appearance correction is always the
cheaper gradient. This prices the conditioning out regardless of how it is
injected, and it bounds what any adapter of this family can achieve without
an objective that pays for actions (\cref{sec:res-decomposition}).

\paragraph{Step-size conditioning through the adapter.} AVID does not
address sampling cost. We extend the composition with a step-size argument
and train it with the shortcut self-consistency objective, so that few-step
behaviour is carried entirely by the adapter while the prior stays frozen.
This is unlike consistency and shortcut models, which retrain the prior. We
show
that this is learnable in the adapter setting, against a matched control, on
two independent statistics (\cref{sec:res-shortcut}), and we derive why the
published velocity-averaging target is exact only on a straight interpolant
and biased on a curved diffusion arc (\cref{sec:curvature}). Few-step
\emph{quality} is directionally better than the no-shortcut control but is
not competitive, and we scope it as future work.

\paragraph{A framework in which the frozen base is a configuration choice.}
We implement the adapter taxonomy behind a single composition interface
spanning diffusion and flow matching, U-Net and DiT. The claim is
demonstrated rather than asserted: the same interface runs across three
backbone families, and we added a new base family to the original AVID
repository as a third branch, running their recipe unchanged
(\cref{sec:res-d1}).

\paragraph{A diagnostic methodology, and the reason it was necessary.} None
of the above is visible in the readouts normally used to monitor such a
model. Loss, gate statistics, sample quality and perceptual metrics are all
blind to whether a conditioned model uses its conditioning. We exhibit a
cell that improves on every quality metric we record while carrying no
recoverable action information (\cref{sec:res-blindness}). We therefore
develop a probe suite that measures the \emph{structure} of the response
rather than its magnitude (\cref{sec:probes}). We regard this as a
contribution in its own right, and we report the case in which our own first
instrument misled us and what was required to correct it
(\cref{sec:integrity}).

\paragraph{What we do not claim.} We do not demonstrate trajectory control.
A paired shuffled-action comparison shows a modest, consistently-signed
effect on how much the arm moves per clip, but supplying the \emph{correct}
action does not reduce prediction error relative to a wrong one on any cell
we measure: what the adapter takes from the action is its magnitude, not its
direction (\cref{sec:res-decomposition}). We do not claim competitive few-step
video generation, a planning or reinforcement-learning contribution, or
state-of-the-art quality on any benchmark. Where a result rests on a single
seed, or on a comparison whose confound we could not remove, we say so at
that point in the text.

\section{Thesis structure}
\label{sec:structure}
